%======================================================================
\section{Search spaces, variants and reconciliation}\label{app:search}
%======================================================================

This appendix lists the search spaces of the trained baselines and the AdvInfoNCE re-implementation, the prior
variants, the reconciliation of the stored results, and the protocol of each related study (Section~\ref{sec:related}).

\subsection{Trained baselines and AdvInfoNCE}\label{sec:grids}

\paragraph{Coat.}
The search space is learning rate in $\{10^{-2},3\cdot10^{-3},10^{-3}\}$, weight decay in $\{10^{-5},10^{-4},10^{-3}\}$
and width in $\{32,64\}$. For LightGCN, LightGCN-pt, r-AdjNorm, NAVIP, DR-LightGCN and SimGCL the depth ranges
over $\{1,2,3\}$. Epochs are limited to 100, the batch size is 256 and the patience is 5. Configurations: LightGCN
54, LightGCN-pt 54, NAVIP 54, DR-LightGCN 54, r-AdjNorm 108 (with $r\in\{0.3,0.7\}$), SimGCL 108 (with weight
$\lambda_{\mathrm{cl}}\in\{0.1,0.5\}$), MACR 36 (with $c\in\{0.1,0.3\}$), PDA 36 (with $\gamma\in\{0.1,0.2\}$),
MF, IPS-MF, DR-MF, BPR-MF, DR-JL and MRDR 18 each, iALS 24 (with $\alpha\in\{3,10,30\}$, regularisation in
$\{1,10,100,1000\}$ and width in $\{32,64\}$) and iALS search space B 18 (with $\alpha\in\{10,30,100\}$ and regularisation
in $\{10,100,300\}$). The iALS runs of the prior family use an epoch limit of 30 on Coat, and the stored iALS run
uses 100. The epoch limits are therefore not the same across the iALS rows.

\paragraph{Yahoo!\,R3.}
The search space is learning rate in $\{10^{-2},3\cdot10^{-3},10^{-3}\}$, weight decay in $\{10^{-5},10^{-4}\}$ (with
$10^{-3}$ added for the matrix models), width 64 and depth 2. For LightGCN, NAVIP and r-AdjNorm the width ranges
over $\{64,128\}$ and the depth over $\{1,2,3\}$. Epochs are limited to 60, the batch size is 1{,}024 and the
patience is 5. Configurations: LightGCN 36, NAVIP 36, r-AdjNorm 72, LightGCN-pt 6, DR-LightGCN 6, SimGCL 12, MACR
18, PDA 18, iALS 12 (with the iALS search space above and width 64), MF, IPS-MF, DR-MF, BPR-MF, DR-JL and MRDR 9 each.

\paragraph{KuaiRand-Pure.}
The search space is learning rate in $\{10^{-2},3\cdot10^{-3},10^{-3}\}$ and weight decay in $\{10^{-5},10^{-4}\}$, with
width 64 and depth 2 for every model. Epochs are limited to 40, the batch size is 8{,}192 and the patience is 3.
Configurations: LightGCN 6, LightGCN-pt 6, NAVIP 6, DR-JL 6, MF 6, DR-MF 6, BPR-MF 6, r-AdjNorm 12, SimGCL 12, MACR
12 and iALS 12. MRDR, IPS-MF, PDA and DR-LightGCN were not run on this dataset.

\paragraph{KuaiRec.}
The search space is learning rate in $\{10^{-2},3\cdot10^{-3},10^{-3}\}$ and weight decay in $\{10^{-5},10^{-4}\}$. For
LightGCN and LightGCN-pt the width ranges over $\{64,128\}$ and the depth over $\{1,2,3\}$, and the other models
use width 64 and depth 2. Epochs are limited to 30, the batch size is 32{,}768, the patience is 3, and each epoch
is capped at 2{,}000{,}000 pairs. Configurations: LightGCN 36, LightGCN-pt 36, r-AdjNorm 12, NAVIP 6,
DR-LightGCN 6, SimGCL 12, MACR 12, iALS 12, MF 6, IPS-MF 6, DR-MF 6, BPR-MF 6, DR-JL 6 and MRDR 6. PDA (12
configurations) and two pointwise negative-sampling variants, LightGCN-pn and DR-LightGCN-pn (24 configurations
each), were also run on KuaiRec and are not in the main comparison.

Across datasets, the graph models were given depth search spaces $\{1,2,3\}$ on Coat, Yahoo!\,R3 and KuaiRec, and only
depth 2 on KuaiRand. Budgets are therefore unequal, and the counts make the difference visible. The selection
rule of Section~\ref{sec:selection} is applied to each search space as it stands, and the sensitivity of the conclusions
to the size of the search space is measured in the results.

\subsection{AdvInfoNCE re-implementation}\label{sec:advsetup}

AdvInfoNCE~\citep{zhang2023advinfonce} is an adversarial contrastive loss. We re-implemented it on the LightGCN
encoder. Our numbers are for this re-implementation and not for the original method, and the official code was
not run. The main simplifications are these.
\begin{itemize}
\item The score is the cosine similarity of the propagated embeddings divided by a temperature $\tau$. Each
positive is contrasted with 64 sampled negatives.
\item The adversary is a pair of embedding tables (users and items, dimension 64) trained with Adam to maximise
the contrastive loss. Its learning rate is $10^{-2}$ on Coat, $10^{-3}$ on Yahoo!\,R3 and KuaiRand, and
$5\cdot10^{-4}$ on KuaiRec. Only this embedding-based adversary is implemented.
\item The adversary budget of the original formulation is replaced by a number of adversarial rounds
$\eta\in\{4,12\}$. A round runs after each epoch from the second epoch on and lasts at least one pass over the
positives and 100 Adam steps. The adversary's KL divergence is logged but not projected onto a ball.
\end{itemize}

The search spaces differ in size, and the budgets are unequal. Coat has 216 configurations: three learning rates, three
weight decays, widths $\{32,64\}$, depths $\{1,2,3\}$, $\tau\in\{0.2,0.75\}$ and $\eta\in\{4,12\}$. Yahoo!\,R3 has
16 configurations with $\tau\in\{0.1,0.28\}$, $\eta\in\{4,12\}$, learning rates $\{3\cdot10^{-3},10^{-3}\}$,
weight decays $\{10^{-5},10^{-4}\}$, width 128 and depth 2. KuaiRand has 8 configurations with
$\tau\in\{0.1,0.28\}$, learning rates $\{10^{-2},3\cdot10^{-3}\}$, weight decay $10^{-4}$, width 64 and depth 2.
KuaiRec has 8 configurations with $\tau\in\{0.2,2.0\}$, learning rates $\{10^{-2},3\cdot10^{-3}\}$, weight decay
$10^{-5}$, width 64 and depth 2. KuaiRand and KuaiRec both use $\eta\in\{4,12\}$. On Coat, AdvInfoNCE thus has 216
configurations against 54 for LightGCN and 108 for r-AdjNorm. A larger search space gives a method more chances to fit
noise in the validation data, so the Coat search space favours AdvInfoNCE and the smaller search spaces work against it. We did
not match budgets for this baseline. The only wider re-tuning in this study is the Coat iALS search space
(Table~\ref{tab:T11}).


\subsection{Prior variants}\label{sec:variants}

The variants below change the additive imputation or the way it enters the score. Each is defined in the code of
the repository. In every variant, the prior is the standardised logit of the variant's imputation, except for
\texttt{pop}, which replaces the offset by a popularity score.
\begin{description}
\item[add.] The imputation of Section~\ref{sec:prior}, fitted with weights $O/\bar p$.
\item[add\_noips.] The same ridge fit with weights $O$ in place of $O/\bar p$, so that the fit uses the exposed
pairs without inverse-propensity weights.
\item[pop.] The imputation, and therefore $Q$, is unchanged. The offset is replaced by the standardised
logarithm of the number of positives of each item in the log, $\log(1+\sum_v O_{vi}Y_{vi})$, broadcast to every
user.
\item[shuf\_item.] The item biases $b_i$ are permuted across items by a fixed permutation (seed 0). The
parameters $\mu$ and $b_u$ are kept, and the imputation and the prior are recomputed from the permuted biases.
\item[user\_only.] The item biases are set to zero, so every item has the same score within a user. The ranking
within a user is then decided by tie-breaking.
\end{description}
In the learners of this analysis, only the prior enters the score. The variant's $Q$ is used only by the
Q-weighted methods, and the prior-anatomy runs do not use them.

The prior-anatomy runs are the following. On Coat, iALS-p and LightGCN-p are run over the ten splits with the
variants add, add\_noips, pop and shuf\_item. On KuaiRec, iALS-p is run over the five splits with the same four
variants on the eight-configuration search space. The prior alone is run for all five variants on all four datasets,
with the floors listed in Section~\ref{sec:prior}. The floor of the Coat LightGCN-p runs is 0.05, and that of the
Coat iALS-p runs is 0.1.


\subsection{Reconciliation of the stored results}\label{sec:reconcile}

Before the analysis, the selection of each (file, method) pair was recomputed from its stored per-epoch records. The
recomputation applies the training-time rule of Section~\ref{sec:selection}, including the patience value recorded
with the records, and its mean test score is compared with the stored mean. Of the 154 (file, method) pairs in the
reconciliation, 151 agree within 0.002 in nDCG. The three that do not agree are DR-LightGCN on KuaiRec (recomputed mean
higher by $0.0033$), DR-LightGCN on Yahoo!\,R3 (higher by $0.0058$) and MACR on Yahoo!\,R3 (higher by $0.0060$). All
three come from two earlier result files that do not record the patience value. The stored selections are reproduced
only under patience values of 3 (KuaiRec) and 5 (Yahoo!\,R3), which the files do not record. The KuaiRec DR-LightGCN
value enters the P1 comparison of Table~\ref{tab:T5}, and the Yahoo!\,R3 DR-LightGCN and MACR values appear in
Table~\ref{tab:T4}. Because the selections cannot be recomputed from the records as stored, we report the stored
values. Every sensitivity number in the results is reported relative to the reconciled selection.


\subsection{Positioning of related studies}\label{sec:positioning}

\begin{table}[htbp]
\centering
\caption{Positioning of related studies by their protocol. Entries follow the protocols as
described in the cited papers; ``not stated'' marks a point that the description does not
settle. Uniform data are the uniformly collected or randomised samples of the benchmark.}
\label{tab:positioning}
\footnotesize
\adjustbox{max width=\textwidth}{%
\begin{tabular}{p{2.6cm}p{1.9cm}p{2.5cm}p{2.0cm}p{2.6cm}p{2.4cm}}
\toprule
Study & Graph backbone & Propensity or imputation & Uniform data in training & Uniform data for validation or selection & Unbiased test benchmarks \\
\midrule
AutoDebias \citep{chen2021autodebias} & no (matrix factorisation) & learned weights and imputation & yes, small sample & yes, small sample & Coat, Yahoo! R3 \\
CausE \citep{bonner2018cause} & no (matrix factorisation) & none & yes, randomised sample & no, validation from logged policy & MovieLens, Netflix \\
CDR \citep{song2023cdr} & no (matrix factorisation) & doubly robust family & no & yes, part of unbiased data & Coat, Yahoo! R3, KuaiRand \\
DCE \citep{kweon2024dce} & no (MF or NCF) & doubly robust loss & no & no, held out from biased training data & Coat, Yahoo! R3, KuaiRec \\
BISER \citep{lee2022biser} & no & self-inverse propensity & no & no, held out from biased training data & Coat, Yahoo! R3 \\
ConAlign \citep{zhang2026conalign} & no (dual tower) & MF-IPS and MF-DR baselines & yes & yes & Coat, Yahoo! R3, KuaiRand-Pure \\
NAVIP \citep{kim2022navip} & yes & inverse propensity on edges & no & pseudo-unbiased popularity sample & Yahoo! R3 and pseudo-unbiased sets \\
AdvDrop \citep{zhang2024advdrop} & yes & IPS and DR baselines only & no, for its own method & not stated & Coat, Yahoo!, KuaiRec \\
PPD \citep{islam2026ppd} & yes, post hoc & none & no & yes on Coat and KuaiRec; biased sample on Yahoo! R3 & Coat, Yahoo! R3, KuaiRec \\
DPAA \citep{dpaa2026} & yes & none & no & validation present, construction not confirmed & Coat, Yahoo! R3, KuaiRec \\
DR-GNN \citep{wang2024drgnn} & yes & none & not stated & not stated & Coat, Yahoo \\
\midrule
This study & yes & IPS-fitted imputation and propensity & no, except Coat shipped propensities & yes & Coat, Yahoo! R3, KuaiRand-Pure, KuaiRec \\
\bottomrule
\end{tabular}}
\end{table}
